% ============================================================
% CHAPTER 3: SYSTEM ANALYSIS AND DESIGN
% ============================================================
\chapter{System Analysis and Design}

This chapter presents a detailed analysis and design of the SafeGuard system. It covers the functional and non-functional requirements, system architecture, UML design diagrams, and module descriptions that collectively define the system's structure and behaviour.

\section{Functional Requirements}

The functional requirements of the SafeGuard system are organized by module and specify the behaviours and capabilities the system must provide.

\subsection{User Authentication Module}

\begin{table}[H]
  \centering
  \caption{Functional Requirements: User Authentication}
  \label{tab:fr-auth}
  \begin{tabularx}{\textwidth}{@{}l l X@{}}
    \toprule
    \textbf{ID} & \textbf{Requirement} & \textbf{Description} \\
    \midrule
    FR-AUTH-01 & User Registration & The system shall allow new users to register by providing first name, last name, email address, phone number, password, and password confirmation. The system shall validate email uniqueness and password strength. \\
    FR-AUTH-02 & User Login & The system shall authenticate users using email address and password. Upon successful authentication, the system shall issue a JWT access token valid for 24 hours. \\
    FR-AUTH-03 & Password Hashing & The system shall hash all passwords using bcrypt before storage. Plain-text passwords shall never be stored in the database. \\
    FR-AUTH-04 & Session Management & The system shall use JWT Bearer tokens for session management. Tokens shall be attached to requests via the Authorization header. \\
    FR-AUTH-05 & Auto-Logout & The system shall automatically log out users when their token expires or when a 401 Unauthorized response is received from the server. \\
    FR-AUTH-06 & Profile Management & The system shall allow authenticated users to view and update their profile information, including name, email, and phone number. \\
    FR-AUTH-07 & Password Change & The system shall allow users to change their password by verifying the current password before accepting the new password. \\
    FR-AUTH-08 & Profile Picture & The system shall allow users to upload a profile picture in JPEG, PNG, or WebP format with a maximum size of 5MB. \\
    \bottomrule
  \end{tabularx}
\end{table}

\subsection{Emergency Contact Module}

\begin{table}[H]
  \centering
  \caption{Functional Requirements: Emergency Contacts}
  \label{tab:fr-contacts}
  \begin{tabularx}{\textwidth}{@{}l l X@{}}
    \toprule
    \textbf{ID} & \textbf{Requirement} & \textbf{Description} \\
    \midrule
    FR-CON-01 & Add Contact & The system shall allow users to add emergency contacts with name, relationship, phone number, email (optional), and priority level. \\
    FR-CON-02 & View Contacts & The system shall display all emergency contacts in a card-based layout with avatar initials and contact details. \\
    FR-CON-03 & Edit Contact & The system shall allow users to modify all fields of an existing emergency contact. \\
    FR-CON-04 & Delete Contact & The system shall allow users to remove an emergency contact with a confirmation dialog. \\
    FR-CON-05 & Search Contacts & The system shall provide search functionality across name, phone, and relationship fields using case-insensitive matching. \\
    FR-CON-06 & Priority Filter & The system shall allow filtering contacts by priority level (low, medium, high, critical). \\
    FR-CON-07 & Priority Levels & The system shall support four priority levels: low, medium, high, and critical, with appropriate visual indicators. \\
    \bottomrule
  \end{tabularx}
\end{table}

\subsection{SOS Alert Module}

\begin{table}[H]
  \centering
  \caption{Functional Requirements: SOS Alert System}
  \label{tab:fr-sos}
  \begin{tabularx}{\textwidth}{@{}l l X@{}}
    \toprule
    \textbf{ID} & \textbf{Requirement} & \textbf{Description} \\
    \midrule
    FR-SOS-01 & SOS Trigger & The system shall provide a prominent SOS button that triggers an emergency alert upon confirmation. \\
    FR-SOS-02 & Geolocation & The system shall automatically capture the user's current latitude and longitude using the browser Geolocation API when an SOS is triggered. \\
    FR-SOS-03 & Alert Types & The system shall support six alert types: manual, automatic, panic, medical, fire, and security. \\
    FR-SOS-04 & Custom Message & The system shall allow users to include a custom message with their SOS alert. \\
    FR-SOS-05 & Alert Status & The system shall track alert status through four states: active, acknowledged, resolved, and cancelled. \\
    FR-SOS-06 & Alert History & The system shall maintain a paginated history of all SOS alerts with filtering by status and search. \\
    FR-SOS-07 & Activity Logging & The system shall log every SOS alert trigger in the activity log with timestamp and details. \\
    \bottomrule
  \end{tabularx}
\end{table}

\subsection{Incident Reporting Module}

\begin{table}[H]
  \centering
  \caption{Functional Requirements: Incident Reporting}
  \label{tab:fr-incident}
  \begin{tabularx}{\textwidth}{@{}l l X@{}}
    \toprule
    \textbf{ID} & \textbf{Requirement} & \textbf{Description} \\
    \midrule
    FR-INC-01 & Report Incident & The system shall allow users to report incidents with type, description, location, severity, and optional image attachment. \\
    FR-INC-02 & Incident ID & The system shall auto-generate a unique incident identifier in the format INC-XXXXXXXX for each report. \\
    FR-INC-03 & Incident Types & The system shall support eight incident types: theft, assault, accident, fire, harassment, vandalism, medical, and other. \\
    FR-INC-04 & Severity Levels & The system shall classify incidents by severity: low, medium, high, and critical. \\
    FR-INC-05 & Status Tracking & The system shall track incident status through: pending, investigating, resolved, and closed. \\
    FR-INC-06 & Geolocation & The system shall support latitude and longitude coordinates for incident location. \\
    FR-INC-07 & Search and Filter & The system shall provide multi-criteria filtering by type, severity, status, and text search. \\
    FR-INC-08 & CRUD Operations & The system shall support creating, viewing, editing, and deleting incident reports. \\
    \bottomrule
  \end{tabularx}
\end{table}

\subsection{Dashboard Module}

\begin{table}[H]
  \centering
  \caption{Functional Requirements: Dashboard}
  \label{tab:fr-dashboard}
  \begin{tabularx}{\textwidth}{@{}l l X@{}}
    \toprule
    \textbf{ID} & \textbf{Requirement} & \textbf{Description} \\
    \midrule
    FR-DASH-01 & Statistics Cards & The system shall display summary statistics: total contacts, total SOS alerts, active SOS alerts, and total incidents. \\
    FR-DASH-02 & Recent Alerts & The system shall display the five most recent SOS alerts with status indicators. \\
    FR-DASH-03 & Recent Incidents & The system shall display the five most recent incidents with severity indicators. \\
    FR-DASH-04 & Activity Log & The system shall display the five most recent activity log entries. \\
    FR-DASH-05 & Quick SOS & The system shall provide a quick-access SOS button on the dashboard. \\
    \bottomrule
  \end{tabularx}
\end{table}

\section{Non-Functional Requirements}

\begin{table}[H]
  \centering
  \caption{Non-Functional Requirements}
  \label{tab:nfr}
  \begin{tabularx}{\textwidth}{@{}l l X@{}}
    \toprule
    \textbf{ID} & \textbf{Category} & \textbf{Requirement} \\
    \midrule
    NFR-01 & Performance & The system shall respond to user interactions within 2 seconds under normal load conditions. \\
    NFR-02 & Security & All passwords shall be hashed using bcrypt with salt. JWT tokens shall expire after 24 hours. \\
    NFR-03 & Data Isolation & Each user shall have access only to their own data. Cross-user data access shall be prevented at the API level. \\
    NFR-04 & Usability & The interface shall be intuitive, requiring no more than 3 clicks to reach any primary feature. \\
    NFR-05 & Accessibility & The system shall provide dark and light themes to accommodate user preferences and accessibility needs. \\
    NFR-06 & Compatibility & The system shall function on modern browsers including Chrome, Firefox, Safari, and Edge. \\
    NFR-07 & Scalability & The decoupled architecture shall support independent scaling of frontend and backend components. \\
    NFR-08 & Reliability & The system shall handle errors gracefully with appropriate user feedback messages. \\
    NFR-09 & Maintainability & The codebase shall follow modular architecture with clear separation of concerns. \\
    NFR-10 & File Security & Uploaded files shall be validated for type and size before storage. Maximum file size: 5MB. \\
    \bottomrule
  \end{tabularx}
\end{table}

\section{System Architecture}

The SafeGuard system follows a classic three-tier architecture with clear separation between the presentation, application, and data tiers. Figure~\ref{fig:architecture} illustrates the high-level system architecture.

\begin{figure}[H]
  \centering
  \begin{tikzpicture}[
    box/.style={rectangle, draw, minimum width=3.5cm, minimum height=1.5cm, text centered, font=\small},
    arrow/.style={-{Stealth[length=3mm]}, thick},
    label/.style={font=\small\itshape},
    tierbox/.style={rectangle, draw, dashed, rounded corners, inner sep=10pt}
  ]

    % Presentation Tier
    \begin{scope}[xshift=0cm]
      \node[tierbox, minimum width=9cm, minimum height=4.5cm, fill=blue!5] (ptier) at (0,0) {};
      \node[font=\small\bfseries, anchor=north] at (ptier.north) {Presentation Tier};
      \node[box, fill=blue!15] (react) at (-1.5, 0.3) {React.js 18\\+ MUI 5};
      \node[box, fill=blue!10] (axios) at (1.5, 0.3) {Axios\\HTTP Client};
      \node[box, fill=blue!8] (router) at (0, -1.2) {React Router\\v6};
      \node[box, fill=blue!12] (ctx) at (0, 1.5) {AuthContext\\ThemeContext};
    \end{scope}

    % Arrow
    \draw[arrow] (0.3, -2.5) -- (0.3, -3.2);
    \draw[arrow] (-0.3, -3.2) -- (-0.3, -2.5);
    \node[label] at (2.2, -2.85) {REST API\\(JSON)};

    % Application Tier
    \begin{scope}[xshift=0cm, yshift=-5.5cm]
      \node[tierbox, minimum width=9cm, minimum height=4.5cm, fill=green!5] (atier) at (0,0) {};
      \node[font=\small\bfseries, anchor=north] at (atier.north) {Application Tier};
      \node[box, fill=green!15] (fastapi) at (-1.5, 0.3) {FastAPI\\Python};
      \node[box, fill=green!10] (auth) at (1.5, 0.3) {JWT Auth\\bcrypt};
      \node[box, fill=green!8] (pydantic) at (0, -1.2) {Pydantic\\Validation};
      \node[box, fill=green!12] (sa) at (0, 1.5) {SQLAlchemy\\ORM};
    \end{scope}

    % Arrow
    \draw[arrow] (0.3, -8) -- (0.3, -8.7);
    \draw[arrow] (-0.3, -8.7) -- (-0.3, -8);
    \node[label] at (2.2, -8.35) {SQL\\Queries};

    % Data Tier
    \begin{scope}[xshift=0cm, yshift=-11cm]
      \node[tierbox, minimum width=9cm, minimum height=2.5cm, fill=orange!5] (dtier) at (0,0) {};
      \node[font=\small\bfseries, anchor=north] at (dtier.north) {Data Tier};
      \node[box, fill=orange!15] (mysql) at (0, 0) {MySQL 8.0\\Database};
    \end{scope}

  \end{tikzpicture}
  \caption{Three-Tier System Architecture}
  \label{fig:architecture}
\end{figure}

\subsection{Presentation Tier}

The presentation tier is implemented using React.js 18, a declarative, component-based JavaScript library for building user interfaces. Material UI 5 provides the component library for consistent and professional UI elements. The tier includes:

\begin{itemize}
  \item \textbf{React Router v6} for client-side routing with route guards
  \item \textbf{Context API} for global state management (authentication and theme)
  \item \textbf{Axios} for HTTP communication with request/response interceptors
  \item \textbf{Emotion} for CSS-in-JS styling
\end{itemize}

\subsection{Application Tier}

The application tier is implemented using Python's FastAPI framework, chosen for its automatic request validation, dependency injection system, and native async support. Key components include:

\begin{itemize}
  \item \textbf{FastAPI Router} for organizing API endpoints by module
  \item \textbf{Pydantic Schemas} for request/response validation and serialization
  \item \textbf{JWT Authentication} using python-jose for token generation and validation
  \item \textbf{bcrypt} for password hashing with salting
  \item \textbf{SQLAlchemy ORM} for database-agnostic data access
\end{itemize}

\subsection{Data Tier}

The data tier uses MySQL 8.0 as the relational database management system, accessed through SQLAlchemy ORM with PyMySQL as the database driver. The schema consists of five normalized tables with proper foreign key relationships and cascade delete rules.

\section{Use Case Diagram}

Figure~\ref{fig:usecase} presents the use case diagram for the SafeGuard system, illustrating the interactions between the actor (User) and the system's primary use cases.

\begin{figure}[H]
  \centering
  \begin{tikzpicture}[
    actor/.style={circle, draw, minimum size=1cm, inner sep=0pt, font=\small},
    usecase/.style={ellipse, draw, minimum width=3cm, minimum height=0.9cm, text centered, font=\small},
    system/.style={rectangle, draw, dashed, rounded corners=15pt, minimum width=10cm, minimum height=9cm},
    line/.style={thick}
  ]

    % System boundary
    \node[system] (sys) at (5.5,0) {};
    \node[font=\small\bfseries] at (5.5, 4.2) {SafeGuard System};

    % Actor
    \node[actor] (user) at (-0.5, 0) {User};
    \node[font=\small] at (-0.5, -1) {\textbf{User}};

    % Use cases
    \node[usecase] (uc1) at (4.5, 3) {Register Account};
    \node[usecase] (uc2) at (4.5, 1.8) {Login / Logout};
    \node[usecase] (uc3) at (4.5, 0.6) {Manage Profile};
    \node[usecase] (uc4) at (4.5, -0.6) {Manage Contacts};
    \node[usecase] (uc5) at (4.5, -1.8) {Trigger SOS Alert};
    \node[usecase] (uc6) at (4.5, -3.0) {Report Incident};
    \node[usecase] (uc7) at (8, 1.8) {View Dashboard};
    \node[usecase] (uc8) at (8, 0.6) {View History};
    \node[usecase] (uc9) at (8, -0.6) {Export Data};

    % Lines
    \draw[line] (user) -- (uc1);
    \draw[line] (user) -- (uc2);
    \draw[line] (user) -- (uc3);
    \draw[line] (user) -- (uc4);
    \draw[line] (user) -- (uc5);
    \draw[line] (user) -- (uc6);
    \draw[line] (user) -- (uc7);
    \draw[line] (user) -- (uc8);
    \draw[line] (user) -- (uc9);

  \end{tikzpicture}
  \caption{Use Case Diagram}
  \label{fig:usecase}
\end{figure}

\section{Activity Diagram}

Figure~\ref{fig:activity-sos} illustrates the activity flow when a user triggers an SOS alert, representing one of the most critical workflows in the system.

\begin{figure}[H]
  \centering
  \begin{tikzpicture}[
    activity/.style={rectangle, rounded corners, draw, minimum width=2.5cm, minimum height=0.8cm, text centered, font=\small},
    decision/.style={diamond, draw, minimum size=1.5cm, text centered, font=\small},
    startend/.style={circle, draw, fill=black, minimum size=0.4cm, inner sep=0pt},
    flow/.style={-{Stealth[length=2mm]}, thick},
    yesno/.style={font=\tiny\itshape}
  ]

    % Start
    \node[startend] (start) at (0, 5) {};

    % Activities
    \node[activity, fill=blue!10] (press) at (0, 4) {Press SOS Button};
    \node[activity, fill=yellow!15] (confirm) at (0, 3) {Confirm Alert};
    \node[activity, fill=green!10] (geo) at (0, 2) {Capture Geolocation};
    \node[activity, fill=orange!10] (check) at (0, 1) {Location Available?};
    \node[activity, fill=red!10] (send) at (0, 0) {Send SOS Alert to Server};
    \node[activity, fill=green!15] (log) at (0, -1) {Log Activity};
    \node[activity, fill=blue!10] (notify) at (0, -2) {Notify Emergency Contacts};
    \node[activity, fill=purple!10] (update) at (0, -3) {Update Alert History};
    \node[startend] (end) at (0, -4) {};

    % Flows
    \draw[flow] (start) -- (press);
    \draw[flow] (press) -- (confirm);
    \draw[flow] (confirm) -- (geo);
    \draw[flow] (geo) -- (check);
    \draw[flow] (check) -- node[yesno, right] {Yes} (send);
    \draw[flow] (check) -- node[yesno, left] {No} (send);
    \draw[flow] (send) -- (log);
    \draw[flow] (log) -- (notify);
    \draw[flow] (notify) -- (update);
    \draw[flow] (update) -- (end);

    % Bypass path from check to send
    \draw[flow] (check.west) -- ++(-1,0) |- node[yesno, near start, below] {Fallback} (send.west);

  \end{tikzpicture}
  \caption{Activity Diagram: SOS Alert Workflow}
  \label{fig:activity-sos}
\end{figure}

\section{Sequence Diagram}

Figure~\ref{fig:sequence-login} depicts the sequence diagram for the user login process, illustrating the interaction between the user, frontend, backend API, and database.

\begin{figure}[H]
  \centering
  \begin{tikzpicture}[
    participant/.style={rectangle, draw, minimum width=2cm, minimum height=0.8cm, text centered, font=\small},
    message/.style={-{Stealth[length=2mm]}, thick},
    selfmsg/.style={thick, decorate, decoration={snake, amplitude=1mm, segment length=3mm}},
    returnmsg/.style={-{Stealth[length=2mm]}, thick, dashed},
    lifeline/.style={dashed, gray},
    xstart=-1.5, gap=3.5
  ]

    % Participants
    \node[participant, fill=blue!10] (user) at (0, 5) {User};
    \node[participant, fill=green!10] (fe) at (3.5, 5) {React Frontend};
    \node[participant, fill=orange!10] (api) at (7, 5) {FastAPI Backend};
    \node[participant, fill=red!10] (db) at (10.5, 5) {MySQL DB};

    % Lifelines
    \draw[lifeline] (user) -- (0, -3);
    \draw[lifeline] (fe) -- (3.5, -3);
    \draw[lifeline] (api) -- (7, -3);
    \draw[lifeline] (db) -- (10.5, -3);

    % Messages
    \draw[message] (0, 4) -- node[above, font=\tiny] {1. Enter credentials} (3.5, 4);
    \draw[message] (3.5, 3.3) -- node[above, font=\tiny] {2. POST /api/login} (7, 3.3);
    \draw[message] (7, 2.6) -- node[above, font=\tiny] {3. SELECT user by email} (10.5, 2.6);
    \draw[returnmsg] (10.5, 2) -- node[below, font=\tiny] {4. User record} (7, 2);
    \draw[message] (7, 1.3) -- node[above, font=\tiny] {5. Verify bcrypt hash} (7, 1.3) edge[selfloop, looseness=10, in=200, out=160, distance=1cm];
    \draw[message] (7, 0.6) -- node[above, font=\tiny] {6. Log activity} (10.5, 0.6);
    \draw[message] (10.5, 0) -- node[above, font=\tiny] {7. INSERT activity log} (10.5, 0) edge[selfloop, looseness=10, in=340, out=20, distance=0.8cm];
    \draw[returnmsg] (7, -0.5) -- node[below, font=\tiny] {8. JWT token + user} (3.5, -0.5);
    \draw[returnmsg] (3.5, -1.2) -- node[below, font=\tiny] {9. Store token, redirect} (0, -1.2);

  \end{tikzpicture}
  \caption{Sequence Diagram: User Login Process}
  \label{fig:sequence-login}
\end{figure}

\section{Class Diagram}

Figure~\ref{fig:class-diagram} presents the class diagram illustrating the data model and relationships between entities in the system.

\begin{figure}[H]
  \centering
  \begin{tikzpicture}[
    classbox/.style={rectangle, draw, minimum width=3cm, minimum height=0.6cm, text centered, font=\small},
    relation/.style={-{Stealth[length=2mm]}, thick},
  ]

    % User
    \node[classbox, fill=blue!10] (user) at (0, 3) {\textbf{User}};
    \node[rectangle, draw, anchor=north west, inner sep=3pt, font=\tiny] (userattrs) at (-1.5, 2.6) {
      \begin{tabular}{l}
        id: int \\
        first\_name: str \\
        last\_name: str \\
        email: str \\
        phone: str \\
        password\_hash: str \\
        role: enum \\
        is\_active: bool
      \end{tabular}
    };

    % EmergencyContact
    \node[classbox, fill=green!10] (contact) at (-5, -1) {\textbf{EmergencyContact}};
    \node[rectangle, draw, anchor=north west, inner sep=3pt, font=\tiny] (conattrs) at (-6.8, -1.4) {
      \begin{tabular}{l}
        id: int \\
        name: str \\
        relationship: str \\
        phone: str \\
        email: str \\
        priority: enum
      \end{tabular}
    };

    % SOSAlert
    \node[classbox, fill=red!10] (sos) at (0, -1) {\textbf{SOSAlert}};
    \node[rectangle, draw, anchor=north west, inner sep=3pt, font=\tiny] (sosattrs) at (-1.5, -1.4) {
      \begin{tabular}{l}
        id: int \\
        latitude: float \\
        longitude: float \\
        address: str \\
        message: str \\
        alert\_type: enum \\
        status: enum
      \end{tabular}
    };

    % IncidentReport
    \node[classbox, fill=orange!10] (incident) at (5, -1) {\textbf{IncidentReport}};
    \node[rectangle, draw, anchor=north west, inner sep=3pt, font=\tiny] (incattrs) at (3.5, -1.4) {
      \begin{tabular}{l}
        id: int \\
        incident\_id: str \\
        incident\_type: enum \\
        description: str \\
        location: str \\
        severity: enum \\
        status: enum \\
        image: str
      \end{tabular}
    };

    % ActivityLog
    \node[classbox, fill=purple!10] (activity) at (0, -5) {\textbf{ActivityLog}};
    \node[rectangle, draw, anchor=north west, inner sep=3pt, font=\tiny] (actattrs) at (-1.5, -5.4) {
      \begin{tabular}{l}
        id: int \\
        action: str \\
        details: str \\
        ip\_address: str \\
        created\_at: datetime
      \end{tabular}
    };

    % Relationships
    \draw[relation] (user.south) -- ++(0,-0.3) -| node[near start, right, font=\tiny] {1} node[near end, left, font=\tiny] {*} (contact.north);
    \draw[relation] (user.south) -- ++(0,-0.5) -| node[near start, right, font=\tiny] {1} node[near end, left, font=\tiny] {*} (sos.north);
    \draw[relation] (user.south) -- ++(0,-0.7) -| node[near start, right, font=\tiny] {1} node[near end, left, font=\tiny] {*} (incident.north);
    \draw[relation] (user.south) -- ++(0,-0.2) |- node[near start, right, font=\tiny] {1} node[near end, below, font=\tiny] {*} (activity.north);

  \end{tikzpicture}
  \caption{Class Diagram: Data Model and Relationships}
  \label{fig:class-diagram}
\end{figure}

\section{Module Description}

The SafeGuard system is organized into the following functional modules:

\subsection{Module 1: Authentication and User Management}

This module handles all aspects of user identity management, including registration, login, profile editing, password changes, and profile picture uploads. It implements JWT-based stateless authentication with 24-hour token expiry. The backend routes are defined in \texttt{app/api/auth.py} and \texttt{app/api/profile.py}, while the frontend state is managed through \texttt{AuthContext.jsx}.

\subsection{Module 2: Emergency Contact Management}

This module provides full CRUD functionality for emergency contacts, supporting name, relationship, phone number, email, and priority classification. Contacts can be searched across multiple fields and filtered by priority level. The UI presents contacts in a card-based layout with avatar initials and colour-coded priority indicators.

\subsection{Module 3: SOS Emergency Alert System}

The SOS module is the core safety feature, providing a prominent one-press button that triggers emergency alerts. Upon confirmation, the system captures geolocation coordinates using the browser's Geolocation API, associates them with the alert, and records the event. Alerts are classified by type (manual, automatic, panic, medical, fire, security) and tracked through status states (active, acknowledged, resolved, cancelled).

\subsection{Module 4: Incident Reporting and Tracking}

This module enables structured incident documentation with eight predefined incident types. Each report generates a unique identifier (INC-XXXXXXXX format) and captures description, location, coordinates, severity, and optional image evidence. Incidents progress through a lifecycle of pending, investigating, resolved, and closed statuses.

\subsection{Module 5: History and Analytics}

This module combines SOS alert history and incident history into a tabbed view with pagination, multi-criteria filtering, and search. It includes CSV export functionality for offline record-keeping. The dashboard aggregates real-time statistics including total contacts, total alerts, active alerts, and total incidents.

\subsection{Module 6: Activity Logging}

A cross-cutting concern implemented across all modules, this system records every significant user action (registration, login, SOS trigger, incident report, profile update, password change) in the activity\_logs table with action type, descriptive details, IP address, and timestamp.
